# Limitations and Open Questions  
## Topological Integrity Gravity (TIG)

---

## 🇬🇧 English Version

### 1. Scope of the Present Work

The current formulation of Topological Integrity Gravity (TIG) should be understood as an **effective, spectrally motivated extension of General Relativity**. It is not claimed to be a complete or fundamental theory of quantum gravity.

---

### 2. Equivalence to f(R) / Starobinsky in the Lowest Limit

In the low-curvature limit, TIG reduces to a curvature-corrected \(f(R)\)-type theory and reproduces the structure of Starobinsky gravity.

Therefore:

> TIG is not distinguishable from \(R+R^2\) gravity at the lowest effective order.

Any physical distinction relies on higher-order spectral contributions and the integrity functional.

---

### 3. Status of the Integrity Functional

The functional

\[
\mathcal{I}[D] = \mathrm{Tr}\log\left(1+\frac{D^2}{\Lambda^2}\right)
\]

is introduced as a natural spectral extension, but:

- its uniqueness is not proven  
- alternative spectral functionals may exist  
- its derivation from a fundamental principle remains open  

---

### 4. Parameter Freedom

The theory contains parameters:

\[
\Lambda,\ \alpha,\ \lambda,\ r_0
\]

While relations between them are established, their absolute values are not fixed by the current framework.

This implies:

- partial degeneracy in parameter space  
- need for observational constraints  
- absence of a predictive UV completion  

---

### 5. Scalar Mode Observability

TIG predicts a scalar curvature mode with mass

\[
m = \frac{\Lambda}{\sqrt{6\alpha}}
\]

However:

- its coupling strength scales as \(g_{\mathrm{eff}} \sim G m^2\)  
- for small masses, the coupling becomes extremely weak  

Thus:

> the scalar mode may lie in observable frequency ranges but remain difficult to detect.

---

### 6. Status of Black Hole Solutions

The regular black-hole metric analyzed in this work:

- is kinematically consistent  
- has finite curvature  
- exhibits a de Sitter core  

However:

> it has not yet been derived as an exact solution of the full TIG field equations.

It should therefore be interpreted as an **effective or motivated solution**, not a rigorously proven one.

---

### 7. Dynamical Stability

While several stability indicators are satisfied:

- \(m^2 > 0\) (no tachyonic scalar)  
- regular core  
- extremal remnant  

a full linear perturbation analysis has not yet been performed.

In particular, open questions include:

- quasinormal mode spectrum  
- inner horizon stability  
- potential mass inflation  
- coupled scalar–metric perturbations  

---

### 8. Role of λ

The parameter \(\lambda\) is introduced as a stability coupling controlling spectral feedback.

However:

- its magnitude is not fixed  
- its direct observability is unclear  
- its independence from \(\alpha\) is not fully established  

Therefore:

> the physical necessity of \(\lambda\) remains to be demonstrated.

---

### 9. Energy Conditions

The effective stress-energy tensor associated with the regular solution:

\[
T^{\mathrm{eff}}_{\mu\nu}
\]

has not yet been fully analyzed with respect to:

- weak energy condition  
- strong energy condition  
- null energy condition  

Violations may occur and require careful interpretation.

---

### 10. Cosmological Implications

While TIG contains \(R^2\)-type terms similar to inflationary models:

- no full cosmological solution is presented  
- no fit to CMB or large-scale structure data is performed  

Thus:

> cosmological viability remains an open question.

---

### 11. Absence of Quantum Completion

TIG is formulated at the level of an effective classical theory.

It does not currently provide:

- a quantization procedure  
- a connection to a fundamental quantum gravity framework  
- a derivation from a UV-complete theory  

---

### 12. Summary of Limitations

The main limitations of the current work are:

1. equivalence to \(f(R)\) at lowest order  
2. non-uniqueness of the integrity functional  
3. partially unconstrained parameter space  
4. weak scalar coupling in observable regimes  
5. absence of exact solution derivation  
6. incomplete stability analysis  
7. unclear physical role of \(\lambda\)  
8. unresolved energy condition behavior  
9. missing cosmological analysis  
10. lack of quantum completion  

---

### 13. Outlook

Future work should focus on:

- deriving exact solutions from the full TIG action  
- performing linear and nonlinear stability analyses  
- constraining parameters via observational data  
- clarifying the physical role of \(\lambda\)  
- extending the framework to cosmology  
- exploring possible quantum interpretations  

---

## 🇩🇪 Deutsche Version

### 1. Geltungsbereich

Die vorliegende Formulierung von TIG ist als **effektive, spektral motivierte Erweiterung der ART** zu verstehen, nicht als vollständige Quantengravitationstheorie.

---

### 2. f(R)-Grenzfall

Im Niedrigkrümmungsbereich reduziert sich TIG auf eine \(f(R)\)-artige Theorie und ist im führenden Term nicht von Starobinsky-Gravitation unterscheidbar.

---

### 3. Integritätsfunktional

Das Funktional ist plausibel, aber:

- nicht eindeutig hergeleitet  
- nicht alternativlos  
- nicht fundamental begründet  

---

### 4. Parameter

\[
\Lambda,\ \alpha,\ \lambda,\ r_0
\]

sind nicht vollständig fixiert.

---

### 5. Skalarer Modus

Der Modus ist prinzipiell messbar, aber:

- oft extrem schwach gekoppelt  
- schwer nachweisbar  

---

### 6. Schwarze Löcher

Die Lösung ist:

- regulär  
- physikalisch motiviert  

aber:

- noch nicht exakt aus TIG abgeleitet  

---

### 7. Stabilität

Nicht vollständig geklärt:

- lineare Störungen  
- QNM  
- innere Horizonte  

---

### 8. Rolle von λ

Physikalisch plausibel, aber:

- nicht fixiert  
- nicht direkt beobachtbar  

---

### 9. Energiebedingungen

Noch nicht vollständig analysiert.

---

### 10. Kosmologie

Noch nicht ausgearbeitet.

---

### 11. Fazit

TIG ist:

> eine konsistente, aber noch unvollständige effektive Theorie.

---

### 12. Ausblick

Zukünftige Arbeiten müssen:

- Stabilität vollständig klären  
- Parameter einschränken  
- Lösungen exakt ableiten  
- Beobachtungen einbeziehen  
